\documentclass[11pt]{article}

% ---------------------------------------------------------------------------
% Project description: Causal structure in the internal representations of a
% data-driven weather model (GraphCast). Audience: advisor + collaborators.
% This document doubles as the seed for the eventual paper.
% ---------------------------------------------------------------------------

\usepackage[margin=1in]{geometry}
\usepackage{amsmath,amssymb}
\usepackage{graphicx}
\usepackage{booktabs}
\usepackage{natbib}
\usepackage{hyperref}
\usepackage{xcolor}
\usepackage{enumitem}

\hypersetup{colorlinks=true, linkcolor=blue!50!black, citecolor=blue!50!black,
            urlcolor=blue!50!black}

% Inline editorial marker for things still to fill in.
\newcommand{\todo}[1]{\textcolor{red}{\textbf{[TODO: #1]}}}
\newcommand{\note}[1]{\textcolor{blue!60!black}{\textit{(#1)}}}

\title{Does a Data-Driven Weather Model Learn Causal Structure?\\
       \large Causal Graphs from Sparse-Autoencoder Features and Mechanisms from
       Parameter Decompositions, Validated on a Synthetic Benchmark}
\author{\todo{authors}}
\date{\today}

\begin{document}
\maketitle

\begin{abstract}
\todo{Write last. One paragraph: weather foundation models forecast well but are
opaque; do their internal representations encode the \emph{causal} structure of the
atmosphere or only correlational shortcuts? Two complementary probes: (i) extend
sparse-autoencoder (SAE) interpretability of GraphCast from single-feature
interventions to systematic \emph{causal discovery} over the graph of features;
(ii) decompose the forecaster's parameters into mechanisms with ablation-defined
causal importance (VPD). Both validated on SAVAR, a synthetic system with a known
ground-truth causal graph, including under temporal subsampling/aliasing and mesh
heterogeneity. Third and most transferable: an answer-key-free selector
(pool-crossed agreement between a read channel and a poke channel) that ranks
candidate feature sets without ground truth, validated out-of-sample across nine
generated worlds and calibrated into a trust dial with a per-domain error bar.}
\end{abstract}

\section{Motivation and Research Question}
\label{sec:intro}

Data-driven weather models---GraphCast~\cite{graphcast}, GenCast, Pangu-Weather,
FourCastNet---now match or beat operational numerical weather prediction on
medium-range skill, at a small fraction of the compute. What they learned in order
to do so is unknown. A physics solver can be read term by term against the
conservation laws it discretizes; a trained graph network cannot. The opacity has
practical costs: forecasts of cyclones, heat waves, and floods feed real decisions,
and a model leaning on training-distribution correlations may fail quietly as the
climate drifts away from that distribution. There is a scientific cost too. If the
model has internalized genuine atmospheric dynamics, its internals amount to a
compressed description of those dynamics that we are currently unable to read.

A model can be highly skillful while having learned only correlational structure.
The question that separates a statistical emulator from a model of the underlying
system is causal:

\begin{quote}
Does GraphCast's internal representation encode the \emph{causal structure} of the
atmosphere, recoverable as a directed graph over interpretable features, or only
the correlational structure sufficient for next-step prediction?
\end{quote}

Accuracy does not settle this. A model at the predictability ceiling may have
organized its computation around the true dependency structure of the atmosphere,
or around shortcuts that predict equally well in distribution; observational
evaluation cannot tell the two apart, which is why interpretability claims of this
kind need interventional or otherwise causal grounding to
generalize~\cite{joshi2026}.

We approach the question along two routes. The first treats sparse-autoencoder
(SAE) features of the model's activations as variables and runs time-series causal
discovery over them. The second decomposes the model's weights into mechanisms
whose importance is defined by ablation (Sec.~\ref{sec:approach}). Neither can be
validated on GraphCast itself, where no ground-truth feature graph exists, so we
validate both on SAVAR, a synthetic spatiotemporal system whose causal graph is
known by construction (Secs.~\ref{sec:savar} and~\ref{sec:ext}).

That last sentence names a problem larger than validation. Both routes begin by
choosing variables---an aggregation of the model's state into a manageable feature
set---and on SAVAR we can pick that aggregation by trying many and keeping
whichever yields the best graph. On GraphCast we cannot, and a method that needs
the answer key to choose its own variables does not transfer. Much of the work
reported here is therefore aimed at a third question: can a candidate feature set
be ranked, and the accuracy of the graph built on it be predicted, using only
quantities computable without ground truth? Section~\ref{sec:selector} answers
yes, by crossing two independent readings of the model---one that \emph{reads}
correlations through a candidate map and one that \emph{pokes} the frozen
forecaster through it---and shows the resulting selector holds up across
overlapping features, rollout-trained forecasters, coupled channels, and the
generator's realism mechanisms, while naming the conditions under which it does
not.

\subsection{Relation to prior work}

MacMillan and Ouellette~\cite{macmillan2025} trained SAEs on GraphCast activations
and found features corresponding to recognizable physics: tropical cyclones,
atmospheric rivers, diurnal and seasonal cycles, precipitation, sea ice, geographic
position. Scaling a single tropical-cyclone feature changed the forecast in a
physically consistent way, so at least some features are causally connected to the
output. What their analysis leaves open, and what they themselves name as the open
problem, is how the features relate to one another. A single-feature intervention
establishes that a feature matters for the forecast; it says nothing about the web
of lagged influences among features, which is where any learned analogue of
atmospheric dynamics would have to live. Recovering that graph is the aim of this
project.

\section{Background}
\label{sec:background}

\subsection{GraphCast}

GraphCast~\cite{graphcast} is an encode--process--decode graph network over the
global atmospheric state. The encoder maps the input fields (atmospheric variables
on pressure levels plus surface fields, on a latitude--longitude grid) onto a
multi-scale icosahedral mesh. A 16-layer message-passing processor, with 512 latent
dimensions per node, propagates information along mesh edges of varying length;
this is where the learned dynamics live. The decoder maps mesh states back to the
grid as a 6-hour forecast, and longer forecasts come from feeding the model its own
output.

Two architectural facts matter later. Rolling the model forward (or teacher-forcing
it through consecutive observed states) supplies a time axis over its internal
features, which the causal-discovery route depends on (Sec.~\ref{sec:temporal}).
And the multi-mesh is heterogeneous: refinement leaves some nodes much better
connected than others. That heterogeneity is a plausible origin for the
position-coding (``grid-locked'') features reported in~\cite{macmillan2025}, and it
is why the synthetic forecaster we build in Sec.~\ref{sec:gnn} uses a heterogeneous
mesh as well.

\subsection{Sparse autoencoders}

An SAE is an overcomplete dictionary trained to reconstruct a layer's activations
from a sparse code, in the hope that the dictionary directions are monosemantic
where the raw neurons are not~\cite{cunningham2023,bricken2023}. The TopK variant
keeps the $K$ largest feature activations per sample. \citet{macmillan2025}
obtained their GraphCast features this way, from processor activations. Their
evaluation stops at feature existence and single-feature interventions; ours
continues to the causal graph among the features.

\subsection{Causal discovery from time series}

We use three discovery methods whose assumptions barely overlap.
PCMCI~\cite{runge2019} is the primary one: a condition-selection stage followed by
momentary-conditional-independence tests, designed for autocorrelated
high-dimensional time series, returning lagged directed edges with significance
levels. PCMCI orients only lagged edges ($\tau \ge 1$); its successor
PCMCI+~\cite{runge2020} also estimates contemporaneous ($\tau = 0$) adjacencies,
which becomes necessary once dynamics act faster than the sampling interval
(Sec.~\ref{sec:subsample}). TSCI~\cite{tsci} comes from the
convergent-cross-mapping tradition and infers causation from reconstructed
attractor geometry. DYNOTEARS~\cite{dynotears} learns a weighted dynamic Bayesian
network by acyclicity-constrained continuous optimization. \note{DYNOTEARS runs in
an isolated venv here due to a numpy/pandas pin; see \texttt{requirements.txt}.}
Where the three agree we gain confidence; where they disagree, the disagreement
itself tells us which assumptions bind.

One identifiability fact shapes several experiments below. Under linear dynamics
with Gaussian noise, the direction of a contemporaneous edge cannot be determined
from observational data: $X \to Y$ and $Y \to X$ induce the same covariance.
Shimizu et al.~\cite{shimizu2006} showed that non-Gaussian noise breaks the tie
(the LiNGAM result: linear, non-Gaussian, acyclic models are fully identifiable,
because higher-order moments carry directional information). The standard ParCorr
test inside PCMCI+ uses second-order statistics only and so cannot exploit this;
Sec.~\ref{sec:subsample} measures what that costs.

\subsection{Parameter decomposition}
\label{sec:pd-background}

SAEs analyze activations. For the weights we use stochastic parameter
decomposition (SPD)~\cite{spd}, which writes each weight matrix as a sum of
rank-one components, $W \approx \sum_c U_{:,c} V_{c,:}$, trained so that the
components sum to $W$, few are needed for any given input, and each is simple
(rank one, single layer). ``Needed'' has an operational definition here. A gate
$g_c(x) \in [0,1]$, learned by randomly masking components and requiring the
masked network to reproduce the original output, estimates how far component $c$
can be attenuated at input $x$ before the computation changes. VPD~\cite{vpd}
hardens the scheme with adversarially chosen masks, so faithfulness must hold
under worst-case ablations, and adds a frequency penalty that pushes each
component to be needed rarely---that is, to specialize.

The gate is the reason we adopted this method. It is an ablation, so causal
importance comes built in; no separate intervention experiment is needed to ask
whether a component is used. Better still for our purposes, when the decomposed
layer is applied at every spatial position with shared weights, as in message
passing, evaluating the gate per position gives each component a spatial usage
map---and the model contains no per-position parameters that could fake this
(Sec.~\ref{sec:vpd}).

\subsection{SAVAR}

SAVAR (Spatially Aggregated VAR)~\cite{savar} is a stochastic spatiotemporal
climate benchmark: a handful of spatial \emph{modes} evolve under a
vector-autoregressive process with a prescribed causal graph and are rendered onto
a grid through fixed mode weight maps. The graph among modes is known by
construction, which is what lets us score a discovery pipeline before trusting it
on GraphCast.

Our default instantiation, all of it part of the given data-generating process:
\begin{itemize}[nosep]
  \item $N=8$ latent modes on a $50\times50$ grid ($L=2500$ pixels).
  \item Mode maps $W$: eight fixed, $L_1$-normalized, non-overlapping Gaussian
        blobs in a $3\times3$ layout. $W$ maps modes to pixels;
        $W^{+}=\mathrm{pinv}(W)$ maps pixels back to modes.
  \item The causal graph $G$: a hand-designed coefficient set
        $\{X_i(t-\tau)\!\to\!X_j(t)\}$, each mode carrying an autoregressive term
        plus cross-edges. It was engineered to contain a hub (X0), a converging
        node with two parents (X3), a pure sink (X7), a weak lag-2 feedback
        (X2$\to$X0), negative edges (X0$\to$X5, X2$\to$X3, X3$\to$X7), and a mix of
        lag-1 and lag-2 structure.
  \item Dynamics and noise: $Z(t)=\sum_\tau \Phi(\tau)Z(t-\tau)+\xi(t)$,
        observations $=W^{+}Z+\text{spatial noise}$, with $\lambda=1$, $D_x=I_N$,
        $D_y=I_L$, giving $\Sigma_y=W^{+}W^{+\top}+I_L$; $T=500$ usable steps after
        a burn-in of $200$. The 100 realisations share $W$ and $G$ and differ only
        in noise seed.
\end{itemize}
The coefficient dictionary $G$ is the ground truth every method in
Sec.~\ref{sec:methodcompare} is scored against. This default generator is linear
and Gaussian; Sec.~\ref{sec:ext} makes it harder in targeted ways (fine temporal
cadence with heterogeneous lags, nonlinear dynamics, skewed innovations, widely
spread mode timescales), each aimed at a question the default cannot pose.

\section{Approach}
\label{sec:approach}

\subsection{Route 1: causal discovery over SAE features}

\begin{enumerate}[nosep]
  \item Run the forecaster over held-out trajectories and record intermediate
        (processor) activations at every step.
  \item Train an SAE on those activations to obtain a sparse feature basis.
  \item For each trajectory, read out each feature's activation at every step,
        giving multivariate time series $\{a_f(t)\}_{f=1}^{F}$.
  \item Run PCMCI+ over the feature series (TSCI and DYNOTEARS serve as
        cross-checks) to recover a lagged directed graph among features.
  \item Score the graph. On SAVAR this means precision, recall, F1, sign accuracy,
        and threshold-free AUROC/AUPRC against the known graph. On GraphCast, where
        no ground truth exists, it means consistency with known physical
        relationships and with interventional edits in the style
        of~\cite{macmillan2025}.
\end{enumerate}

Agreement among the three discovery methods on SAVAR, where we can score them,
calibrates how much to trust each of them later on GraphCast, where we cannot.

\subsection{Route 2: parameter decomposition}

The second route applies VPD (Sec.~\ref{sec:pd-background}) to the forecaster's
shared message-passing MLPs. Each weight matrix becomes a set of rank-one
components with a per-node gate; averaging the gate over inputs and reshaping to
the grid gives every component a spatial usage map. Since the components themselves
are shared across all nodes, exactly as in GraphCast's processor, any position
specificity in those maps had to emerge through the gate.

The two routes answer different questions and inherit different weaknesses. SAE
features live in activation space and come with a time axis, so they support
claims about dynamics, but their causal role must be established after the fact
and the time axis brings aliasing problems (below). VPD components live in weight
space, their importance is an ablation by construction, and no time axis is
involved; what they cannot directly show is how the represented quantities evolve.

\subsection{Do SAE features support temporal causal claims?}
\label{sec:temporal}

An SAE is trained per activation vector, so a feature is an atemporal descriptor
of a single state, and it is fair to ask whether causal discovery, which is about
time, is even well posed on such features. It is, because the time axis never
comes from the SAE. It comes from the trajectory: the SAE supplies the variables,
the forecaster's evolution across steps supplies the dynamics, and reading a
feature's activation across steps gives an ordinary time series.

The argument comes with caveats, which we treat as real limitations and return to
in Sec.~\ref{sec:open}. A feature has to keep its meaning from one step to the
next, or its time series is incoherent. On GraphCast the collection regime
matters: free-running rollouts accumulate error and drift off-distribution, so the
series is non-stationary and conditional-independence tests are confounded,
whereas teacher-forced trajectories (consecutive single steps on observed states)
stay approximately stationary. We therefore collect on teacher-forced trajectories
unless closed-loop behaviour is itself under study. Discovery on internal features
also recovers the causal structure of the \emph{model's} computation, which need
not coincide with the atmosphere's; for the question ``what did the model learn?''
the model's structure is the right target, but the claim should be stated as such.
And the 6-hour step bounds the resolvable lag: anything faster aliases into
contemporaneous dependence. Section~\ref{sec:subsample} measures this effect and
what PCMCI+ recovers of it.

\section{Validation on SAVAR: Default Setting (Completed Work)}
\label{sec:savar}

This section establishes, on the default SAVAR instantiation, that the
SAE$\rightarrow$causal-discovery pipeline recovers true causal structure, and that
the recovered structure connects to what the forecaster actually computes.

\note{All numbers below are from the current repository. Verify against the latest
result files before finalizing.}

\subsection{The forecaster reaches the predictability ceiling}

A spatiotemporal CNN trained to forecast the SAVAR observation field attains
validation RMSE $1.072$, against an oracle floor of $1.061$ and a persistence
baseline of $1.48$. The forecaster is near-optimal, which makes its representations
a fair place to ask what a well-performing model learns.

\subsection{Discovery on the latent modes recovers the ground-truth graph}

PCMCI run on the true latent mode series recovers the ground-truth graph with
precision $71.5\%$, recall $98.4\%$, F1 $82.5\%$, and $100\%$ sign accuracy. So the
discovery method is adequate for SAVAR dynamics before any SAE enters the picture,
and this F1 of $0.825$ becomes the ceiling the rest of the pipeline is measured
against.

\subsection{Discovery on SAE features}

The same discovery run on SAE-feature time series gives precision $0.629$, recall
$0.789$, F1 $0.695$, sign accuracy $100\%$. Causal structure survives the SAE
bottleneck with a moderate loss (F1 $0.825 \to 0.695$); the features are imperfect
proxies for the modes, but causally informative ones.

\subsection{Interpretable features require per-mode SAEs}

A single SAE trained on all modes jointly fails, and the failure has a clean
explanation: the mode-weighted activations have degenerate PCA structure, with PC0
carrying ${\approx}86\%$ of the variance as a shared global-activity direction, so
every mode collapses onto the same feature. Per-mode SAEs repair this. On the
diurnal forecaster $8/8$ modes align, $5/8$ strongly, $2/8$ monosemantically; on
the default forecaster $7/8$ align, none strongly. Part of the gap is a hard
limit---ridge-regression ceilings show some modes are simply not linearly
recoverable from these activations above $r=0.5$, no matter the SAE. Even so,
polysemanticity is the rule in every configuration we tried, and
Sec.~\ref{sec:open} discusses what that means for defining causal variables.

\subsection{Grid-locked features are decodable but causally inert}
\label{sec:gridlock-cnn}

A spatial SAE on per-pixel activations finds many \emph{grid-locked}
(position-coding) features. Ablation shows the forecast does not use them: zeroing
the grid-locked directions removes $23\%$ of the activation norm and moves RMSE by
$-0.02\%$. The controls make the contrast sharp. Six \emph{content} features
carrying $1.8\%$ of the norm cost $+10.1\%$ RMSE when ablated; a random set of the
same size costs $+3.0\%$. Position information sits in the representation without
being exploited by the computation. GraphCast has geographic-coding features of its
own~\cite{macmillan2025}, so this distinction will matter there.

\subsection{Forecast importance tracks content and causal out-degree}

Per-feature ablation ranks the content features as the main contributors to
forecast skill, with the grid-locked features exactly inert, and the ranking holds
out-of-sample. (The in-sample correlation between content-$R^2$ and importance,
$+0.43$, vanishes out-of-sample, so we make no quantitative claim there.)

Against the recovered causal graph, forecast importance follows \emph{driving}:
out-degree correlates with importance at Spearman $+0.78$, out-strength at $+0.76$,
descendant count at $+0.68$, while the in-flavored centralities (PageRank,
in-degree, ancestor count) come out negative. Amplitude confounds this (variance
alone correlates with importance at ${\approx}+0.7$), and partialling it out leaves
total-degree ($+0.56$) and a weakened out-degree ($+0.36$). Centrality does predict
importance, but only once amplitude is controlled.

\subsection{Method comparison: PCMCI vs.\ TSCI vs.\ DYNOTEARS}
\label{sec:methodcompare}

Table~\ref{tab:methods} scores all three discoverers against the same ground-truth
graph, on the clean latents $Z$ and on the SAE-feature series. The threshold-free
AUROC/AUPRC on the summary graph are the numbers to compare across methods
(matching the CausalDynamics-style scoring in \texttt{eval\_common.py}); the
lag-resolved precision/recall/F1 depend on each method's detection threshold
($\alpha$ for PCMCI, $|w|$ for DYNOTEARS) and are only loosely comparable.

\begin{table}[h]
\centering
\small
\begin{tabular}{llccccc}
\toprule
Method & Run on & Prec. & Recall & F1 & AUROC & AUPRC \\
\midrule
PCMCI     & latent   & 0.715 & 0.984 & 0.825 & ---   & ---   \\
DYNOTEARS & latent   & 0.839 & 0.995 & \textbf{0.908} & \textbf{1.000} & 0.998 \\
TSCI      & latent   & 0.214 & 0.499 & 0.299 & 0.500 & 0.277 \\
\midrule
PCMCI     & features & 0.629 & 0.789 & \textbf{0.695} & \textbf{0.923} & 0.863 \\
DYNOTEARS & features & 0.262 & 0.669 & 0.376 & 0.725 & 0.623 \\
TSCI      & features & 0.215 & 0.502 & 0.301 & 0.502 & 0.262 \\
\bottomrule
\end{tabular}
\caption{Causal-graph recovery on SAVAR (default dataset, 100 realisations). Sign
accuracy is $100\%$ for the signed methods (PCMCI-ParCorr, DYNOTEARS) on both rows.
PCMCI-latent AUROC was not computed by the original \texttt{run\_pcmci.py}.}
\label{tab:methods}
\end{table}

DYNOTEARS is nearly perfect on the clean latents (AUROC $1.000$) and falls apart
under the SAE bottleneck (AUROC $0.725$); PCMCI loses the least (AUROC $0.923$).
Since activation-derived series are precisely the GraphCast setting, PCMCI is the
method to lean on there. TSCI sits at chance throughout (AUROC ${\approx}0.50$),
which is what one should expect: the default SAVAR is a linear low-dimensional VAR,
and there is no nonlinear attractor geometry for cross-mapping to exploit. The
chance-level result delimits where state-space methods apply; it is not an
implementation failure.

\subsection{Summary}

On a system with known ground truth, the pipeline works end-to-end. Causal
structure comes through the SAE bottleneck at a moderate cost, the recovered
topology relates to forecast importance, and ablation cleanly separates
representational artifacts from content the forecast depends on. What these
results cannot yet address---temporal aliasing, mesh-induced position coding,
identifiability of contemporaneous edges---is the subject of the next section.

\section{Extensions: Aliasing, Mesh Heterogeneity, and Parameter Decomposition}
\label{sec:ext}

The validation above used the default linear-Gaussian generator and a CNN
forecaster, and several GraphCast-relevant questions cannot even be posed in that
setting. Discovery ran at the cadence the dynamics are defined on, so nothing ever
aliased. A translation-equivariant CNN has no per-position parameters and no
privileged locations, so grid-locked features cannot arise in it, whatever the
data. And linear-Gaussian dynamics leave contemporaneous directions unidentifiable
in principle while pinning the forecaster's representation to a nearly low-rank
subspace. The experiments in this section relax each of these constraints.

\subsection{A harder generator}
\label{sec:finecadence}

The fine-cadence generator (\texttt{generate\_finecadence.py}; 100 realisations,
$T=2400$) keeps the eight-mode edge set of the default graph and changes the rest.
Cross-edges get heterogeneous lags $\ell \in \{1,2,3,4,6\}$ in fine-step units. The
dynamics gain a saturating autoregression and a bilinear interaction term
($\alpha=0.5$, $\beta=0.15$). Innovations become skew-normal (shape $\alpha=4$,
standardized to unit variance, so VAR stationarity is untouched). The skewness is
there for a reason: it is the LiNGAM condition under which contemporaneous
directions become identifiable in principle (Sec.~\ref{sec:background}), so this
generator can tell us whether our discovery methods actually cash in on that
identifiability.

\subsection{Subsampling and aliased edges: PCMCI+ vs.\ PCMCI}
\label{sec:subsample}

Observing the fine-cadence system at stride $s$ turns a ground-truth lag $\ell$
into a coarse lag $\lfloor \ell/s \rfloor$, so every edge with $\ell < s$ aliases
to a contemporaneous ($\tau=0$) dependency. Plain PCMCI with $\tau_{\min}=1$ has no
way to represent such an edge; its recall on aliased edges is zero before it sees
any data. Table~\ref{tab:subsample} shows the sweep over strides for PCMCI+
($\tau_{\min}=0$) and PCMCI ($\tau_{\min}=1$) on the latent modes (40 realisations,
$\alpha=0.05$).

\begin{table}[h]
\centering
\small
\begin{tabular}{cccccc}
\toprule
$s$ & $T$ & GT edges (lag/contemp.) & PCMCI+ contemp.\ F1 & PCMCI+ overall F1 &
PCMCI overall F1 \\
\midrule
1 & 2400 & 12 / 0 & ---   & 0.776 & 0.593 \\
2 & 1200 &  9 / 3 & 0.627 & 0.587 & 0.457 \\
3 &  800 &  6 / 6 & 0.491 & 0.444 & 0.308 \\
4 &  600 &  4 / 8 & 0.376 & 0.346 & 0.192 \\
\bottomrule
\end{tabular}
\caption{Subsampling sweep on the fine-cadence dataset. ``GT edges'' counts
ground-truth edges that remain lagged vs.\ alias to $\tau=0$ at stride $s$.
Contemporaneous F1 is computed on the undirected adjacency. At $s=1$ the PCMCI+
lagged F1 is $0.825$, matching the default-dataset ceiling of
Sec.~\ref{sec:savar}.}
\label{tab:subsample}
\end{table}

PCMCI+ recovers a substantial fraction of the aliased edges (contemporaneous F1
$0.38$--$0.63$) while plain PCMCI's overall F1 collapses with aliasing, $0.593$
down to $0.192$. Whenever dynamics might act faster than the sampling
interval---and at 6-hour steps on the real atmosphere they certainly do---PCMCI+
is the tool. Its remaining weakness is orientation. Most recovered $\tau=0$ edges
come back unoriented (orientation rate $0.08$--$0.30$), even though the skewed
innovations make the directions identifiable in principle. ParCorr is a Gaussian
test statistic; the directional information lives in higher moments it never
looks at. Bolting a LiNGAM-style orientation step onto PCMCI+'s adjacencies is
the obvious follow-up and has not been done yet (Sec.~\ref{sec:open}).

\subsection{A GraphCast-faithful forecaster}
\label{sec:gnn}

For grid-locked features to be possible at all, the forecaster needs privileged
locations. We replaced the CNN with a message-passing GNN
(\texttt{train/gnn\_forecaster.py}) built like GraphCast's processor in miniature:
a mesh over the $50\times50$ grid with local 8-neighbour edges everywhere plus
long-range edges among a lattice of hub nodes, node degrees running from 3 to 32,
and message-passing MLPs shared across all nodes. The mesh is the only source of
positional asymmetry. The SAVAR generator supplies none: the mode layout is
content, tied to the causal graph, and the observation noise is spatially
homogeneous. Any hub-tied structure in the trained model is therefore attributable
to the mesh.

An early version of this model also had a learnable per-node embedding, added on
the reasoning that a CNN lacks per-position parameters and the GNN should have
them. We removed it. GraphCast's processor has no such lookup table, and giving
position a parametric home of its own would have manufactured the grid-lock
phenomenon we were trying to observe emerging from the mesh. All results below use
the model without it.

The trained GNN reaches validation RMSE $0.4295$ on the fine-cadence one-step
task, against $0.4591$ for predicting the grand mean: $R^2 = 0.125$. This is a
weak-signal task by construction (the one-step ceiling at the mode level is
$R^2 \approx 0.20$; Sec.~\ref{sec:homog}), so throughout this section we read the
structure of the interpretability results and put little weight on absolute effect
sizes.

\subsection{Per-mode SAEs on the GNN}
\label{sec:gnn-sae}

We repeated the per-mode SAE pipeline of Sec.~\ref{sec:savar} on the GNN's node
states after the final message-passing layer (mode-weighted pooling, eight
independent TopK SAEs, ridge ceilings measured out-of-fold). The picture is the
CNN's, only sharper. Alignment drops from $7/8$ modes to $2/8$ (X1 and X2), still
with no strong alignment and no monosemantic feature anywhere. The global-activity
direction that dominated the CNN's pooled activations at ${\approx}86\%$ of
variance dominates the GNN's at $88$--$92\%$, and the linear decoding ceilings
shrink accordingly ($0.36$--$0.50$, versus $0.36$--$0.59$ for the CNN). The SAEs
themselves are not the bottleneck; they still reach $66$--$81\%$ of ceiling. No
feature is mode-specific (specificity $\le +0.07$), and modes X1 and X6 literally
share their best feature. Message passing over the heterogeneous mesh concentrates
node state onto global system activity even more aggressively than convolution
does, and mode-specific structure never surfaces in the representation. The
parameter-level analysis below says something similar about the hub structure.

\subsection{Parameter decomposition of the message-passing MLPs}
\label{sec:vpd}

We applied VPD to the frozen GNN: both linear maps of all four message-passing
MLPs, $C=64$ rank-one components each, per-node gates, stochastic plus adversarial
(PGD) masks, frequency-minimality on. Reconstruction converged: masked-forward
MSE to the frozen model's output of ${\approx}2\times10^{-3}$ under stochastic
masks, ${\approx}3.5\times10^{-2}$ under worst-case masks, faithfulness residual
${\approx}8\times10^{-4}$. The decomposition reproduces the forecaster even under
worst-case node-level ablations, which is the property that licenses reading the
gates causally.

Each component's gate map (mean gate per node, reshaped to $50\times50$) is
summarized by its \emph{hub contrast} (mean gate on hub nodes minus background),
its \emph{mode contrast} (mean gate on the eight mode blobs minus background), and
an ablation effect (rise in reconstruction error when the component is fully
masked). Table~\ref{tab:vpd} gives per-layer means.

\begin{table}[h]
\centering
\small
\begin{tabular}{lccccc}
\toprule
Layer & Hub contrast & Mode contrast & Abl.\ mean ($\times10^{-3}$) &
Abl.\ max ($\times10^{-3}$) & Mode-dominant \\
\midrule
0, MLP in   & \textbf{0.479} & 0.369 & 1.86 & 17.3 & 19/64 \\
0, MLP out  & \textbf{0.258} & 0.113 & 0.62 &  2.1 & 18/64 \\
1, MLP in   & \textbf{0.441} & 0.327 & 3.08 & 37.1 & 21/64 \\
1, MLP out  & \textbf{0.222} & 0.149 & 0.20 &  0.6 & 26/64 \\
2, MLP in   & 0.151 & \textbf{0.340} & 2.40 & 25.1 & 50/64 \\
2, MLP out  & 0.006 & \textbf{0.227} & 0.10 &  1.2 & 48/64 \\
3, MLP in   & 0.118 & \textbf{0.365} & 1.24 &  8.2 & 48/64 \\
3, MLP out  & 0.041 & 0.049 & 0.03 &  0.9 & 41/64 \\
\bottomrule
\end{tabular}
\caption{VPD gate structure per decomposed matrix (64 components each).
``Mode-dominant'' counts components whose gate is at least as mode-locked as
hub-locked. Bold marks the larger of the two contrasts.}
\label{tab:vpd}
\end{table}

Depth organizes the gates. In layers 0--1 they concentrate on the mesh hubs: hub
contrast exceeds mode contrast, and only 19--26 of 64 components are more mode-
than hub-locked. By layers 2--3 the pattern has inverted, with 48--50 of 64
components tracking the SAVAR mode blobs. Early mechanisms lock to the mesh, late
mechanisms to content. Nearly all of this structure sits in the first sublayer of
each MLP, whose gate input includes the neighbour-aggregated state; that is enough
to tell a hub from an ordinary cell without any positional encoding.

The gates also settle the grid-lock question, this time from inside a single
object. Across early-layer components, hub contrast correlates negatively with
ablation effect: $-0.48$ at layer 0, $-0.37$ at layer 1. The most hub-locked
mechanisms are the ones the forecast depends on least. Section~\ref{sec:gridlock-cnn}
reached the same conclusion for the CNN's SAE features, but needed a separate
ablation experiment to get there; here the gate is ablation-defined, so
``position-locked'' and ``causally inert'' are read off the same map. No time
axis is involved either, so the aliasing issues of Sec.~\ref{sec:subsample} never
arise. For grid-lock specifically, the parameter route is simply the better
instrument.

Caveats: the ablation effects are small in absolute terms everywhere (means of
$1$--$3\times10^{-3}$), as expected against a weak-signal target, so the depth
gradient and the hub--ablation correlation are claims about structure, and we make
no claim about magnitudes. The components are also far from specialized: every
component's gate covers all eight mode blobs almost equally (coefficient of
variation ${\approx}0.05$). Why nothing specializes is the next question.

\subsection{Why nothing specializes}
\label{sec:homog}

The GNN's margin over the grand-mean predictor is thin ($R^2 = 0.125$), and its
behaviour is essentially damped persistence: correlation $0.78$ with the last
input frame, effective gain ${\approx}0.27$. The generator explains this. The
one-step mode-forecast ceiling is only $R^2 \approx 0.20$, and the eight modes are
dynamically near-interchangeable (variance ratio $1.1$, self-coupling
coefficients all within $[0.30, 0.55]$, cross-coupling worth a mere $0.03$ of
$R^2$). For such a system the optimal one-step forecaster is close to one shared
shrinkage toward the mean, applied identically everywhere. Mode-specific
circuitry would buy nothing, and the model built none; the SAE features of
Sec.~\ref{sec:gnn-sae} and the VPD components of Sec.~\ref{sec:vpd} are
homogeneous because the task is.

This account makes a testable prediction: give the modes genuinely different
dynamics and specialization should appear. A modified generator
(\texttt{generate\_hetdynamics.py}) widens the self-coupling range to
$[0.15, 0.92]$, a ${\sim}23\times$ spread in autocorrelation timescale, with an
equal-variance control to separate timescale heterogeneity from amplitude. The
three-way comparison (baseline, heterogeneous dynamics, equal-variance control)
through the full GNN and SAE pipeline is running at the time of writing.
\todo{Add heterogeneous-dynamics results when the runs complete.}

\section{Choosing a Mode Map Without an Answer Key}
\label{sec:selector}

Everything so far was scored against ground truth. GraphCast offers none, and the
gap this opens is not a detail of evaluation but a hole in the middle of the
method. The pipeline of Sec.~\ref{sec:approach} begins by choosing a set of
variables---an aggregation map from the model's high-dimensional state to a
handful of candidate features---and that choice dominates everything downstream.
On SAVAR we can simply try many maps and keep the one with the best graph. On
GraphCast we cannot, so a recipe that needs the answer key to pick its own
variables is not a recipe at all.

This section builds a selector that ranks candidate maps using only quantities
computable without ground truth, and calibrates how much its verdict can be
trusted. It is the part of the project that has moved most since the results
above, and the part the GraphCast application depends on.

\subsection{The candidate pool and the two channels}
\label{sec:px}

The construction needs a pool of candidate aggregation maps spanning a wide
quality range, and two independent ways of reading a causal graph through each.

The pool (\texttt{sae/discover\_modes.py}) holds thirteen candidates: unsupervised
discovery operators (varimax-PCA, $k$-means and DMD clustering, each run on raw
pixels and on model activations), the oracle map, and deliberately corrupted
variants---blurred, shifted, merged, split, diagonal---that manufacture a
controlled quality axis. Scored against truth the pool spans F1 $0.00$ to $0.89$,
which is what a selector test requires.

Two channels then read a graph through each candidate map $\hat W$. The
\emph{int} channel pools the observation field through $\hat W$ and runs PCMCI+
over the resulting series, giving $\hat G_{\text{int}}$: this is discovery by
\emph{reading} correlations. The \emph{dyn} channel injects impulses into the
frozen forecaster along $\mathrm{pinv}(\hat W)$, reads the teacher-forced response
back out on the rows of $\hat W$, and thresholds against a pixel-permutation null,
giving $\hat G_{\text{dyn}}$: discovery by \emph{poking} the model. Neither uses
the true graph, the true modes, or the true footprints anywhere.

An earlier attempt at unsupervised selection scored candidates by the internal
consistency of a single channel (agreement between aggregate-level and
pixel-level independence verdicts). It reached Spearman $+0.52$ against true
quality and failed for a diagnosable reason: a map that destroys all signal has no
dependencies at either level and is therefore vacuously consistent. Consistency is
a trustworthy \emph{screen}---its zeros always mean a dead map---but not a
\emph{ranker}. That failure is what motivated crossing two channels instead.

\subsection{Same-map agreement fails, in two informative ways}

The natural two-channel statistic is the agreement between the two graphs read
through the \emph{same} map, $F_1(\hat G_{\text{int}}(A), \hat G_{\text{dyn}}(A))$.
It does not work: Spearman $+0.375$ against true quality, far below the
pre-registered bar of $0.8$. Table~\ref{tab:samew} shows why. The two failures are
structural, and both carry to GraphCast.

\begin{table}[h]
\centering
\small
\begin{tabular}{lccc}
\toprule
Candidate & True F1 & Same-$\hat W$ agreement & Pool-crossed (PX) \\
\midrule
blur      & 0.918 & 0.556 & 0.453 \\
oracle    & 0.879 & 0.621 & 0.447 \\
shift5    & 0.878 & \textbf{0.154} & 0.490 \\
vmax\_act & 0.835 & 0.298 & 0.424 \\
dmd\_act  & 0.632 & 0.118 & 0.351 \\
coarse4   & 0.230 & \textbf{0.727} & 0.149 \\
diag8     & 0.000 & 0.000 & 0.000 \\
\bottomrule
\end{tabular}
\caption{Selected candidates from the parent pool ($13$ total). Same-$\hat W$
agreement is maximized by the worst usable map (\texttt{coarse4}) and minimized by
one of the best (\texttt{shift5}); pool-crossing repairs both.}
\label{tab:samew}
\end{table}

First, \textbf{merged maps self-agree on the wrong graph}. The candidate
\texttt{coarse4}, which merges mode pairs and recovers almost nothing
(true F1 $0.230$), attains the pool's \emph{highest} same-map agreement
($0.727$). Both channels inherit the same marginalization from the shared $\hat W$
and therefore make the same mistake in the same place. Agreement measures
consistency of error, not correctness, whenever the error is shared.

Second, \textbf{the write channel is geometry-sensitive where the read channel is
not}. The candidate \texttt{shift5} has footprints in the wrong locations but
behaviourally coherent pooled series; it reads a nearly correct graph
(true F1 $0.878$) while its dyn channel returns almost nothing (agreement
$0.154$). Reading through $\hat W$ needs only that the pooled series be
behaviourally coherent; \emph{writing} through $\mathrm{pinv}(\hat W)$ needs the
injected perturbation to land on the true supports. This asymmetry is worth
stating on its own: perturbation-based validation implicitly certifies footprint
\emph{geometry}, not just behaviour, and a map can pass one test while failing the
other.

\subsection{Pool-crossed agreement}

Both failures come from the two channels sharing a map. Crossing the pool removes
the shared term. Define
\[
  \mathrm{PX}(A) \;=\; \frac{1}{|\mathcal{P}|-1}\sum_{B \neq A}
  F_1\!\left(\hat G_{\text{int}}(A),\, \hat G_{\text{dyn}}(B)\right),
\]
the mean agreement between candidate $A$'s read graph and every \emph{other}
candidate's poke graph. The dyn pool acts as a $\hat W$-independent reference, so
per-map channel error averages out while genuine map quality survives. The
direction matters: the int channel carries the quality signal, and the transposed
statistic (each candidate's dyn graph against the int pool) scores only $+0.288$.

On the parent rung PX reaches Spearman $+0.950$ (Pearson $+0.99$) against true
graph quality, against $+0.375$ for same-map agreement. Good maps sit at
$0.40$--$0.49$, merged or weak maps at $\le 0.17$, dead maps at exactly $0$---a
usable selector with a clean gap.

One caveat belongs here rather than in a footnote: the aggregation rule
(mean over the pool) was fixed \emph{after} the two confounds above were
diagnosed, on $13$ candidates. The parent number is therefore in-sample, and the
rungs below exist to test it out-of-sample.

\subsection{A scoring convention that had to be corrected}

Diagnosing \texttt{shift5} exposed an error in the ground truth itself, not in the
selector. Under the original convention---matching discovered variables to true
modes by footprint cosine---\texttt{shift5} scored F1 $0.000$. But its pooled
series correlate $0.979$ with the true modes; matching variables to modes by
series correlation instead gives F1 $0.835$. The map was good and the score was
wrong. The same correction lifts \texttt{dmd\_act} from $0.177$ to $0.553$.

All scoring in this section therefore uses behaviour-based matching. The
methodological point generalizes beyond SAVAR: \emph{footprint geometry is the
wrong success criterion for an aggregation map; behavioural coherence of the
pooled series is the right one}. On GraphCast, where features have no
ground-truth spatial support to be compared against, only the behavioural
criterion is available at all.

\subsection{Out-of-sample rungs}

Each rung in Table~\ref{tab:rungs} changes one property of the world or the
forecaster, retrains where needed, and reruns the battery unchanged. The rule was
pre-registered from the parent's saved battery and applied verbatim, with no
tuning; the bar is Spearman $\ge 0.8$.

\begin{table}[h]
\centering
\small
\begin{tabular}{llcc}
\toprule
Rung & What changes & PX Spearman & Verdict \\
\midrule
parent & --- (in-sample reference)            & $+0.950$ & --- \\
R1 & overlapping footprints (max cosine $0.20$) & $+0.946$ & \textbf{pass} \\
R3 & $C=2$ channels with cross-channel edges    & $+0.943$ & \textbf{pass} \\
R5 & rollout-trained forecaster (BPTT, horizons 1--8) & $+0.934$ & \textbf{pass} \\
R4b & $N=24$ modes, $N$ unknown a priori        & $+0.790$ & miss \\
\bottomrule
\end{tabular}
\caption{Pool-crossed selector across pre-registered rungs. Each is a full
generate--train--discover--calibrate chain; the dyn channel was verified live in
every one, so each number is a real test rather than an inapplicable channel.}
\label{tab:rungs}
\end{table}

The three passes matter for different reasons. \textbf{R1} removes the
disjoint-footprint assumption that every earlier SAVAR variant relied on, and it
is the sharpest test of the two-channel idea because overlap degrades the
internals channel specifically: discovery from activations collapses (\texttt{vmax\_act}
graph F1 $0.819 \to 0.185$) while the pixel side stays healthy ($0.938$). The
selector nonetheless ranks the pool correctly. \textbf{R5} fine-tunes the
forecaster on multi-step rollout, which is GraphCast's actual training curriculum;
it measurably moved the model's internals (\texttt{vmax\_act} truth-F1
$0.485 \to 0.556$) while leaving the pixel side identical, and the selector
survived the shift. \textbf{R3} adds coupled channels and cross-channel edges, and
passes even though its activation-side discovery is degenerate---the multivariate
GNN carries one hidden vector per \emph{spatial} node, so co-located channels
receive identical pooled activations. PX works there because it crosses the clean
pixel-side read channel against the live poke channel.

\textbf{R4b} is the honest miss. Tripling the mode count to $N=24$ with $N$
unknown, on a pre-registered resolution-homogeneous pool of ten candidates, PX
reaches $+0.790$ and misses its bar by $0.010$. It is reported as a miss and was
not re-run. Two things temper it. The Pearson correlation is $+0.975$, so the
calibration line is intact and only the ranking is noisy---with ten candidates a
single adjacent swap costs about $0.02$ Spearman. And the discovery side scaled
cleanly: unknown-$N$ community detection recovered $\hat N = 24$ and matched
$23/24$ modes at cosine $0.966$, while the fixed-rank varimax operators that
worked at $N=8$ collapsed entirely. That contrast is itself a result for the
GraphCast setting, where the number of features is large and not known in advance.

\subsection{Mechanism robustness}

A separate matrix (Table~\ref{tab:robust}) stresses the selector against the
realism mechanisms of the generator---non-Gaussian innovations, non-linear
dynamics, seasonality, and all of them combined---each as a full rung.

\begin{table}[h]
\centering
\small
\begin{tabular}{llcc}
\toprule
Mechanism & PX Spearman & Pearson & Verdict \\
\midrule
Non-Gaussian innovations (skew-normal) & $+0.909$ & $+0.946$ & \textbf{pass} \\
All mechanisms combined                & $+0.869$ & $+0.951$ & \textbf{pass} \\
Non-linearity (bilinear, $\alpha = 0.5$) & $+0.624$ & $+0.972$ & rank miss \\
Seasonality, raw series                & $+0.132$ & $+0.333$ & miss \\
Seasonality, deseasonalized            & $+0.769$ & $+0.789$ & near miss \\
\bottomrule
\end{tabular}
\caption{Selector under generator realism mechanisms. Same-$\hat W$ agreement
stays weak throughout ($+0.18$ to $+0.42$), so pool-crossing continues to carry
the signal.}
\label{tab:robust}
\end{table}

The combined world passing at $+0.869$ is the headline: with seasonality,
saturating and bilinear non-linearity, and skewed innovations all active at once,
the selector still ranks the pool. The two misses are more useful than the passes.

Under strong non-linearity PX becomes a screen rather than a ranker. It reliably
kills the dead maps, but eight candidates sit within $0.017$ of each other in PX
while spanning $0.29$ of true F1, and ranking inside that cluster is noise
(excluding two degenerate candidates the Spearman falls to $+0.382$). What
survives is the decision that actually matters: the top-pick costs only $0.057$ F1
relative to the best available map, and Pearson $+0.972$ says the calibration line
is undamaged. A related check confirms the linear conditional-independence test is
not the problem---under bilinear dynamics ParCorr is edge-recall-preserving, at a
cost of about four spurious edges, while a nonparametric test buys F1 $0.815 \to
0.957$ for a $50{,}000\times$ increase in compute.

Seasonality is the one mechanism that breaks the recipe outright on raw series
($+0.132$) and is repaired by preprocessing ($+0.769$). Shared periodic forcing
induces dependence between co-periodic modes that the conditional-independence
tests read as causal structure. \textbf{Deseasonalization is therefore a required
preprocessing step, not an option}---a concrete and slightly uncomfortable
finding for GraphCast, whose most prominent published SAE features include the
diurnal and seasonal cycles themselves~\cite{macmillan2025}.

\subsection{The trust dial and its per-domain calibration}
\label{sec:dial}

Ranking candidates is half of what is needed. The other half is turning an
observed PX into a predicted accuracy with an error bar---a dial that says how
much to believe a graph recovered from a model with no answer key.

Pooling all candidate points across the original five rungs ($n=52$) gives
\[
  \text{accuracy} \;\approx\; 1.48 \cdot \mathrm{PX} + 0.13,
  \qquad \text{Spearman } +0.905,\;\; R^2 = 0.797,\;\;
  \text{residual s.d.\ } 0.131 .
\]
The four mechanism worlds above were generated after this fit and never entered
it, so they constitute a genuine out-of-sample test ($n=52$ new points). The
result splits cleanly, and the split is the main finding of this subsection.

\emph{Ranking transfers}: Spearman $+0.810$, Pearson $+0.926$ on the unseen
points. \emph{Absolute calibration does not}: only $8$ of $52$ fall inside the old
$\pm 0.13$ band, with a systematic per-world offset of $+0.12$ to $+0.15$---the
old line \emph{under}-predicts accuracy, which is at least the safe direction. A
refit pooled over all nine worlds ($n=104$) gives
$\text{accuracy} \approx 1.76 \cdot \mathrm{PX} + 0.13$ with residual $\pm 0.16$.

The obvious suspect for the offset was a confound in our own procedure: the later
worlds switched conditional-independence test, so the shift might have been ours
rather than the world's. Rerunning R1 identically under the new test settles it.
PX and true accuracy move together (mean $+0.046$ and $+0.049$), which slides
points \emph{along} the calibration line rather than above it; the residual goes
from $-0.048$ to $-0.067$, against the $+0.136$ the mechanism worlds show. The
test choice contributes essentially nothing. The offset is a property of the
data-generating domain, which strengthens rather than weakens the conclusion: what
must be recalibrated per domain is the intercept, not the instrument.

The resulting claim, stated at the strength the evidence supports: \textbf{PX is a
rank-reliable selector (out-of-sample Spearman $+0.81$) whose absolute
accuracy calibration is per-domain, to about $\pm 0.16$ F1}. It is not a universal
constant, and quoting it as one would be the easiest way to oversell this work.

Three preconditions emerged from the rungs, and all three are detectable without
ground truth---which is what makes them usable rather than merely known. The
candidate pool must be resolution-homogeneous (the R4 lesson: crossing graphs of
different sizes muddies the pairwise F1, and pool resolution spread is directly
measurable). The dyn channel must be live (an atmosphere-regime rung produced a
dead poke channel, which makes the selector inapplicable rather than wrong; the
number of detected dynamical edges reveals this immediately). And the input series
must be deseasonalized.

\subsection{A by-product: disagreement localizes model deficiencies}

The two channels were built to be crossed, but their \emph{disagreements} turn out
to be informative on their own. On maps whose poke channel is live, the true edges
that appear in the read graph but not in the poke graph are almost exactly the
couplings the forecaster failed to internalize: three specific edges, $15/16$
detections, with no non-ground-truth disagreements. Pooled over all good maps
including those with handicapped poke channels, $62\%$ of disagreements localize a
genuine model deficiency against a $25\%$ base rate.

``In the graph you read but not the graph you poke'' is thus a candidate
model-evaluation product in its own right: it points at where an emulator has
learned the correlational structure of a coupling without implementing the
mechanism. It survives losing the oracle map, so it is available in the GraphCast
setting, where it would name the atmospheric couplings the model predicts through
but does not implement.

\section{Open Questions and Risks}
\label{sec:open}

\paragraph{Feature identity across time.} The time axis comes from the trajectory
(Sec.~\ref{sec:temporal}), and the whole construction assumes a feature means the
same thing at step $t$ and step $t{+}1$. If a feature drifts in meaning, every
causal edge built on its time series is suspect. We still owe this an explicit
stability check on GraphCast rollout data.

\paragraph{Model dynamics vs.\ atmospheric dynamics.} Discovery on internal
features recovers the causal structure of the model's computation. Whether that
matches the atmosphere's causal structure is a separate empirical question. The
interpretability framing (``what did the model learn?'') legitimately targets the
model's structure, but every claim should say which of the two it is about.

\paragraph{No ground truth in GraphCast.} SAVAR lets us score recovered graphs;
GraphCast never will. Validation there has to lean on known physical relationships
(the diurnal cycle driving near-surface temperature, say) and on interventional
consistency in the style of~\cite{macmillan2025}: an edge $A \to B$ produced by
discovery should survive an actual intervention on $A$.

\paragraph{Polysemanticity.} In every SAVAR variant tested, the best-aligned
features remain polysemantic, and a polysemantic feature is an ambiguous causal
variable. Possible responses include stricter monosemanticity criteria, per-region
SAEs, or grouping features before discovery. The heterogeneous-dynamics experiment
(Sec.~\ref{sec:homog}) will show whether specialization appears once the task
rewards it; if it does not, the grouping problem moves to the center of the
project.

\paragraph{Orientation of contemporaneous edges.} The aliasing risk itself is now
quantified: Sec.~\ref{sec:subsample} shows plain PCMCI collapsing under
subsampling while PCMCI+ recovers the aliased adjacencies. Time priority still
orients every lagged edge, and a time-unrolled lagged graph is a DAG, so feedback
loops such as X2$\to$X0 stay identifiable. The unsolved part is direction at
$\tau=0$: ParCorr orients only $8$--$30\%$ of the contemporaneous edges it finds,
despite innovations skewed enough to make the directions identifiable in
principle. A LiNGAM-style orientation step on top of the PCMCI+ adjacencies is the
natural fix and is not yet implemented. On GraphCast at 6-hour cadence, sub-step
atmospheric pathways will alias, and the sweep in Table~\ref{tab:subsample} is a
quantitative preview of what that costs.

\paragraph{Weak-signal regimes.} The fine-cadence one-step task leaves the
forecaster barely above climatology, capping the effect size any ablation can
show (Sec.~\ref{sec:vpd}). Conclusions there rest on which components matter and
where, never on how much. GraphCast at 6 hours is far above climatology, so this
is a SAVAR-specific worry for now, but it will return at long lead times as
rollout skill decays.

\paragraph{Overlapping modes.} Every SAVAR variant so far uses spatially disjoint
mode maps, so the pixel-to-mode projection is clean. Overlapping modes (closer
centers, wider spatial profiles) would degrade the projection and dilute the SAE
features. The knob exists in the generator and has not been turned.

\paragraph{Method disagreement on GraphCast.} If PCMCI+, TSCI, and DYNOTEARS
disagree on GraphCast features, some adjudication rule is needed, and it should be
calibrated on SAVAR. Section~\ref{sec:methodcompare} is the start of one: on
activation-derived series, weight PCMCI+ most, and treat edges only DYNOTEARS
reports with caution.

\section{Contributions and Plan}
\label{sec:contributions}

\subsection{Intended contributions}

\begin{enumerate}[nosep]
  \item Extend SAE interpretability of GraphCast from single-feature
        interventions~\cite{macmillan2025} to causal-graph discovery over the
        feature set.
  \item Validate the pipeline on SAVAR's known ground truth: SAE features support
        causal-structure recovery (latent F1 $0.825$ vs.\ feature F1 $0.695$), and
        the recovered topology relates to forecast importance.
  \item Method guidance for the no-ground-truth setting, from comparing PCMCI+,
        TSCI, and DYNOTEARS on activation-derived series, including under temporal
        subsampling. The sweep in Sec.~\ref{sec:subsample} says when PCMCI+ is
        required and locates the remaining orientation gap.
  \item Separate decodable-but-inert representational artifacts (grid-locked and
        geographic-coding features) from content the forecast depends on, via
        causal ablation on the SAE side and ablation-defined gates on the
        parameter side.
  \item A parameter-space route: VPD on a message-passing forecaster yields
        spatial mechanism-usage maps, these show a depth gradient from mesh-locked
        to content-locked mechanisms, and the grid-lock-is-epiphenomenal result
        comes out of a single ablation-defined object (Sec.~\ref{sec:vpd}). To our
        knowledge this is the first application of the SPD family to a GNN.
  \item \textbf{An answer-key-free selector for causal variables}
        (Sec.~\ref{sec:selector}). Pool-crossed agreement between a read channel
        (discovery through a candidate map) and a poke channel (impulse response
        through the frozen model) ranks candidate feature sets with no ground
        truth, at Spearman $+0.93$ to $+0.95$ on three out-of-sample rungs
        ($+0.95$ in-sample), where same-map agreement reaches only $+0.38$. We
        give the two structural reasons the naive statistic fails---shared-map
        error and the read/write geometry asymmetry---and both carry beyond
        SAVAR.
  \item \textbf{A calibrated trust dial with an honest scope limit.} Observed
        agreement maps to predicted graph accuracy; the ranking transfers
        out-of-sample (Spearman $+0.81$ on four unseen worlds) while the absolute
        calibration is per-domain to ${\pm}0.16$ F1. A dedicated rerun rules out
        our own conditional-independence-test switch as the cause of the offset.
  \item \textbf{Preconditions that are detectable without ground truth}: a
        resolution-homogeneous candidate pool, a live perturbation channel, and
        deseasonalized inputs. Seasonality breaks the recipe on raw series
        ($+0.13$) and is repaired by preprocessing ($+0.77$), which matters
        directly for GraphCast, whose published SAE features prominently include
        the diurnal and seasonal cycles.
  \item An attribution discipline for negative results: the missing mode
        specialization was traced to the generator's dynamic homogeneity, and the
        experiment testing that explanation split---confirming it for the SAE
        representation and refuting it for the parameter decomposition, which
        relocates that failure from the data to VPD's objective
        (Sec.~\ref{sec:homog}).
\end{enumerate}

\subsection{Milestones}

\todo{Confirm/adjust with the advisor.}
\begin{enumerate}[nosep]
  \item SAVAR validation writeup (Sec.~\ref{sec:savar}) --- done.
  \item PCMCI/TSCI/DYNOTEARS comparison table --- done.
  \item Fine-cadence generator; subsampling sweep; switch to PCMCI+ --- done.
  \item Heterogeneous-mesh GNN forecaster; per-mode SAEs on it --- done.
  \item VPD decomposition of the GNN's message-passing MLPs --- done.
  \item Heterogeneous-dynamics experiment (specialization test) --- done; split
        result (Sec.~\ref{sec:homog}).
  \item Answer-key-free selector; pre-registered rungs for overlap, rollout
        training, multivariate channels, and scale --- done
        (Sec.~\ref{sec:selector}).
  \item Mechanism-robustness matrix and trust-dial calibration --- done.
  \item Scale rung at a fair pool: $+0.790$ against a bar of $0.8$. A larger
        pre-registered pool at $N=24$ is the outstanding follow-up.
  \item LiNGAM-style orientation of contemporaneous edges --- not started.
  \item GraphCast: activation collection on teacher-forced trajectories; SAE
        training on the target model itself.
  \item GraphCast: feature-stability check; unsupervised feature grouping
        (Ising/Leiden) before discovery; deseasonalization of all derived series
        before conditional-independence testing.
  \item GraphCast: causal discovery on features; selector and trust dial
        recalibrated in-domain; physical-plausibility and
        interventional-consistency evaluation.
  \item Write-up.
\end{enumerate}


\bibliographystyle{plainnat}
\bibliography{refs}

\end{document}
